%% Part VII -- Social Utopias (comic pages 50 to 53).
%% The only section of the book set as running prose rather than panels: four
%% pages of essay that ask what we would point the technology at if we got to
%% choose. The frames follow the essay's own order, and the last one is the
%% question the book ends on.
%% Fragment only: no preamble, \input by ai_for_all.tex.

\sectionsubtitle{pages 50 to 53}
\section{Social Utopias}

% Page 51. The time comparison is the section's opening move and the best
% single antidote to the word revolution: spears for 300,000 years, song
% recommendations for under ten.
\begin{frame}{AI today is not Skynet}
  \panel{page 50, the section opener}
  \begin{comicquote}{51}
    For over 300,000 years we've been hunting with spears; for less than ten
    years we've been looking for our next song with the help of AI.
  \end{comicquote}
  \slidesources{Schneider2019}
\end{frame}

% Page 52, and the conditional the whole utopia hangs on. Quoted in full
% because the second half, reduce not increase, is the part that gets
% dropped when this argument is repeated from memory.
\begin{frame}{It depends on the training data}
  \panel{page 52, the application procedure}
  \begin{comicquote}{52}
    If we train AIs on discriminatory data, they will make discriminatory
    decisions and reduce, not increase, the opportunities, and resources of
    the underprivileged.
  \end{comicquote}
  \slidesources{Schneider2019}
\end{frame}

% Page 52 again, the motorway example: an AI that documents every
% inhabitant's involvement transparently and proposes solutions. The claim
% quoted is the modest one the authors actually make for it.
\begin{frame}{Raising the question, not settling it}
  \panel{page 52, the motorway through the city}
  \begin{comicquote}{52}
    An AI may not be able to resolve all injustices, but it can raise
    questions about how we want to live.
  \end{comicquote}
  \slidesources{Schneider2019}
\end{frame}

% Page 53. The authors turn on their own utopia in one paragraph, which is
% the move that makes the section worth teaching. Ask the room to answer the
% third question before moving on.
\begin{frame}{Or a creepy surveillance state}
  \panel{page 53, the same platform, read twice}
  \begin{comicquote}{53}
    Or is this vision making us create a creepy surveillance state? Why
    should the authorities always be well-meaning? How could we make sure
    they are?
  \end{comicquote}
  \slidesources{Schneider2019}
\end{frame}

% The book's last words, and the deck's assignment. Tbc. is the authors'.
\begin{frame}{What do you think?}
  \panel{page 53, Tbc.}
  \begin{comicquote}{53}
    But the first step has been made: What can AI help us with? Where do we
    have to be careful? What do you think? Let your voice be heard and your
    point of view be seen.
  \end{comicquote}
  \slidesources{Schneider2019}
\end{frame}

% Page 55, the book's own reading list, carried over whole. It is the shelf
% the essay was written off, and it is the only frame in the part that is not
% a quotation, so it takes a tabular rather than a comicquote.
\begin{frame}{Further reading, the book's own}
  \panel{page 55, the reading list}
  \begin{sidebody}
    \scriptsize
    \renewcommand{\arraystretch}{1.15}%
    \begin{tabular}{@{}p{0.30\linewidth}p{0.62\linewidth}@{}}
      \toprule
      \multicolumn{1}{@{}l}{\alert{Authors}} & \multicolumn{1}{l@{}}{\alert{Title}}\\
      \midrule
      Agrawal, Gans, Goldfarb & Prediction Machines: The Simple Economics of Artificial Intelligence, 2018\\
      Evans                   & Broad Band: The Untold Story of the Women Who Made the Internet, 2018\\
      Chollet, Allaire        & Deep Learning with R, 2018\\
      Lee                     & AI Superpowers: China, Silicon Valley, and the New World Order, 2018\\
      Bostrom                 & Superintelligence: Paths, Dangers, Strategies, 2014\\
      Zuboff                  & The Age of Surveillance Capitalism, 2018\\
      Russell, Norvig         & Artificial Intelligence: A Modern Approach, 2010\\
      Daum                    & Das Kapital sind wir: Zur Kritik der digitalen \"Okonomie, 2017\\
      \bottomrule
    \end{tabular}
  \end{sidebody}
  \slidesources{Schneider2019}
\end{frame}
